La principal aportación de este Trabajo de Fin de Grado es un framework
\emph{full-stack} reutilizable que reduce drásticamente el código que hay que
escribir y mantener para exponer, validar, comunicar y editar una entidad de
datos cuando el modelo está en evolución. Concretamente, el trabajo aporta:

\begin{itemize}
  \item Un \textbf{contrato de comunicación uniforme} (envoltorio de
  respuesta, catálogo de errores y descubrimiento) que desacopla el frontend
  del dominio del backend.
  \item Una \textbf{cadena de generación automática} que convierte el modelo
  ORM en esquema enriquecido y este, sin intervención manual, en formularios
  renderizables.
  \item Una \textbf{capa cliente genérica tipada} en Angular que convive con
  el código heredado sin romperlo, facilitando una migración progresiva.
  \item Un \textbf{entorno de despliegue reproducible} multiplataforma con
  Docker, validado en sistemas operativos y arquitecturas distintas.
\end{itemize}

La repercusión esperada es doble. En el plano técnico, el \gls{itc} dispone de
una base que acorta el tiempo de arranque y el coste de mantenimiento de sus
futuros proyectos de gestión de datos. En el plano
socio-económico, el primer beneficiario es el \gls{jbvc}, cuya labor de
conservación e investigación de la flora canaria se apoyará en aplicaciones
construidas más rápido, con menos errores y con mayor calidad y adaptabilidad.
